\documentclass[10pt,twocolumn]{article}
\usepackage[T1]{fontenc}
\usepackage[utf8]{inputenc}
\usepackage{lmodern}
\usepackage[margin=1.8cm,top=2.4cm,bottom=2cm,columnsep=0.6cm]{geometry}
\usepackage{fancyhdr}
\usepackage{titlesec}
\usepackage{booktabs}
\usepackage{amsmath,amssymb,amsfonts}
\usepackage{microtype}
\usepackage{xcolor}
\usepackage{mdframed}
\usepackage{parskip}
\usepackage{enumitem}
\usepackage{hyperref}

\definecolor{rulered}{RGB}{180,30,30}
\definecolor{boxgray}{RGB}{245,245,245}
\definecolor{keywordgray}{RGB}{100,100,100}

\pagestyle{fancy}
\fancyhf{}
\fancyhead[L]{\textsc{\small Claude Mag}}
\fancyhead[C]{\small\textit{ICML 2026 Special Edition: Uncertainty \& Conformal Prediction}}
\fancyhead[R]{\small 2026-07-08}
\fancyfoot[C]{\small\thepage}
\renewcommand{\headrulewidth}{0.8pt}

\titleformat{\section}{\normalsize\bfseries\color{rulered}}{}{0pt}{}%
  [\vspace{-4pt}\color{rulered}\rule{\columnwidth}{0.4pt}]
\titlespacing{\section}{0pt}{8pt}{4pt}

\hypersetup{colorlinks=true,urlcolor=rulered,linkcolor=black}

\begin{document}
\setlength{\parindent}{0pt}
\setlength{\parskip}{4pt}

{\small\color{keywordgray}\textbf{Keywords:}
  conformal risk control \textbullet{}
  generative models \textbullet{}
  large language models \textbullet{}
  memorization}\par\vspace{4pt}

{\LARGE\bfseries\raggedright
  Conf-Gen: Conformal Uncertainty Quantification for Generative
  Models}\par\vspace{4pt}

{\large\itshape\raggedright
  Extending Formal Risk Guarantees Beyond Supervised Prediction into
  LLMs, Image Generators, and Agents}\par\vspace{6pt}

{\small\textbf{By} Gabriel Loaiza-Ganem, Kevin Zhang, Wei Cui, Marc Law
  \& Kin Kwan Leung
  (Layer 6 AI, TD Bank Group)}\par\vspace{2pt}

{\small\color{keywordgray}arXiv:2605.28920 \textbullet{}
  ICML 2026 Poster}
\par\vspace{2pt}\noindent{\color{rulered}\rule{\columnwidth}{0.6pt}}\vspace{6pt}

Conformal risk control (CRC) gives formal uncertainty guarantees for
supervised predictors---but a diffusion sample or an LLM's answer has
no ground-truth label to calibrate against. \textbf{Conf-Gen} relaxes
CRC's assumptions just enough to reach generative systems directly,
producing certified guarantees such as ``this image is not a memorized
copy,'' ``this dialogue system has asked enough clarifying questions,''
and ``this agent's output is correct''---unifying and generalizing
prior one-off attempts to bring conformal guarantees to LLMs.

\begin{mdframed}[backgroundcolor=boxgray,linecolor=rulered,linewidth=1pt,
    innerleftmargin=6pt,innerrightmargin=6pt,
    innertopmargin=6pt,innerbottommargin=6pt]
\textbf{\small AT A GLANCE}\par\vspace{4pt}
\begin{description}[leftmargin=0pt,labelsep=4pt,itemsep=2pt,parsep=0pt,
    font=\small\bfseries]
  \item[Problem:] Conformal risk control assumes labeled prediction
    pairs; generative and agentic systems have no natural label to
    calibrate against.
  \item[Why it matters:] As generative and agentic systems displace
    classical supervised predictors, the uncertainty-quantification
    toolkit built for supervised learning risks becoming obsolete.
  \item[Method:] Relax CRC's exchangeable-labeled-pair requirement to
    an exchangeable-generation-instance requirement, calibrating
    task-specific risk functionals directly on generated outputs.
  \item[Result:] A single calibration recipe transfers across image
    generation, dialogue, and agentic tasks without redesign.
  \item[Evidence:] Demonstrations on non-memorization guarantees for
    image generators, clarification-sufficiency guarantees for
    dialogue systems, and correctness guarantees for agent outputs.
\end{description}
\end{mdframed}

\section{The Big Picture}

Conformal prediction and its extension, conformal risk control, are the
most practically trusted formal uncertainty guarantees in machine
learning---but both were built for supervised prediction, where a
labeled calibration pair $(x, y)$ is available. Modern AI progress is
increasingly driven by unsupervised generative models: LLMs, image and
video generators, and autonomous agents, none of which produce a
prediction against a classical ground-truth label. Conf-Gen asks which
of CRC's assumptions are actually load-bearing, and shows that relaxing
just the labeled-pair requirement is enough to bring the same
finite-sample guarantee machinery to generative tasks.

\section{Why Existing Approaches Fall Short}

Prior attempts to bring conformal guarantees to LLMs typically bolted a
supervised-style calibration step onto a generative pipeline in a
task-specific, one-off way---for instance treating a single decoding
step as if it were a classification label. These approaches do not
generalize: each new generative task (image synthesis, dialogue,
agentic tool use) required its own bespoke reformulation, with no
shared theoretical foundation guaranteeing correctness across tasks.

\section{Core Mechanism}

Conf-Gen replaces the labeled-pair requirement with task-specific risk
functionals defined directly over generated outputs---distance to the
training set as a memorization score, a count of clarifying questions
asked, or a verifier's judgment as a correctness score. Calibration
proceeds exactly as in CRC: find the smallest threshold for which the
empirical risk, corrected for finite-sample slack, does not exceed the
target level. What changes is the object being calibrated: a single
generative process's own outputs, rather than a labeled prediction
task.

\section{Technical Formulation}

Let $R(\lambda) = \mathbb{E}[\ell(\lambda, X)]$ be the population risk
as a function of threshold $\lambda$, for a monotone bounded loss
$\ell$. Calibration selects
\[
\hat\lambda = \inf\Big\{\lambda : \hat{R}_n(\lambda) + c_n(\delta)
  \le \alpha\Big\},
\]
where $\hat{R}_n$ is the empirical risk over $n$ calibration instances
and $c_n(\delta)$ is a finite-sample correction term ensuring
$R(\hat\lambda) \le \alpha$ with probability at least $1-\delta$.
Conf-Gen's relaxation replaces the requirement that $(X,Y)$ be an
exchangeable labeled pair with the weaker requirement that the sequence
of \emph{generation instances} (and their associated risk scores) be
exchangeable---sufficient to preserve the guarantee while accommodating
memorization scores, clarification-sufficiency scores, and
agent-correctness scores as the risk functional.

\section{Learning or Inference Procedure}

\begin{enumerate}[itemsep=2pt,parsep=0pt]
  \item Define a task-specific, monotone risk functional over generated
    outputs (memorization distance, clarification count, verifier
    score).
  \item Collect an exchangeable sequence of generation instances and
    compute the empirical risk at a grid of thresholds $\lambda$.
  \item Calibrate $\hat\lambda$ as the smallest threshold whose
    finite-sample-corrected empirical risk meets the target $\alpha$.
  \item At deployment, apply the calibrated threshold to new generation
    instances, inheriting the risk guarantee.
\end{enumerate}

\section{What the Guarantee Says}

For any monotone bounded risk functional and any exchangeable sequence
of generation instances, $R(\hat\lambda) \le \alpha$ holds with
probability at least $1-\delta$ in finite samples---matching the
guarantee strength of classical CRC, despite the absence of a
conventional labeled calibration set.

\section{Experimental Findings}

The paper demonstrates Conf-Gen across three qualitatively different
generative settings: certifying that image generator outputs are not
memorized copies of training data, certifying that a conversational AI
has asked sufficient clarifying questions before committing to an
answer, and certifying the correctness of an AI agent's output---showing
the same calibration recipe transfers across modalities and task types
without redesign.

\section{Ablations and Interpretation}

\textbf{Unification:} Conf-Gen recovers prior ad hoc LLM conformal
methods as special cases, showing they were instances of a single
underlying relaxation rather than unrelated tricks.

\textbf{Generality:} The framework extends to tasks with no natural
notion of a held-out label at all, such as certifying properties of an
agent's multi-step trajectory.

\textbf{Practical takeaway:} Task designers need only supply a monotone
risk functional over generation instances; the calibration and
guarantee machinery is otherwise unchanged from classical CRC.

\section*{Reference}

Gabriel Loaiza-Ganem, Kevin Zhang, Wei Cui, Marc Law, Kin Kwan Leung.
\textbf{Conf-Gen: Conformal Uncertainty Quantification for Generative
Models}.
arXiv:2605.28920, May 2026.
\url{https://arxiv.org/abs/2605.28920}

\end{document}
